% GPT-6 Astra Pro; /root/fixed_surface; 22 November 2026 UTC.
% Effort not exposed. Separate complete proof for independent review.
\section[HD-F1-S32: arbitrary ambient perfect tensors]{HD-F1-S32: every actual ambient perfect tensor meeting the full S14 character has a nonzero invariant trace kernel}
\label{sec:f1-norm-fourfold-perfect-tensor}

Retain the actual norm quotient algebra
$E=i_*\mathcal O_Y$ on the geometric open of S29.
Let $N\in\operatorname{Perf}(P)$ be any actual perfect
complex, and form the actual derived tensor
\[
             F=E\otimes^{\mathbf L}_{\mathcal O_P}N.
 \tag{F1S32.1}
\]
Its differentials and extension maps are retained.
Write $r=\operatorname{rk}N\in\mathbf Z$ for the
virtual rank, constant because $P$ is connected.
The scalar projector $p_{\rm ns}$ is the actual
quotient-unit projector (F1S31.7).
The graded partial trace constructed below satisfies
\[
 p_{\rm ns}\operatorname{ptr}_N
           \operatorname{ev}_F(\kappa+\eta)
        =-r\,\operatorname{ob}_E(\kappa),
 \qquad
 \kappa\in\mathfrak D_P,\quad
 \eta\in H^2(P,\mathcal O_P).
 \tag{F1S32.2}
\]
Here $\operatorname{ob}_E$ is precisely the S24
first-jet obstruction, and $\operatorname{ev}$ is
the BF characteristic action with the explicitly
calibrated sign (F1S31.12a).

Every such actual tensor satisfying the full S14
requirement
\[
 e^B\operatorname{ch}(F)=\lambda t+Q(\beta),\qquad
 B\in H^{1,1}(P,\mathbf Q),\quad
 0\ne\lambda t\in W_K(P),\quad Q\in\mathbf Q[x]
 \tag{F1S32.3}
\]
has $r\ne0$. For every specified coherent lift of
a faithful finite point subgroup, its invariant
full degree-two trace kernel has dimension at least
$16-r_C\ge4$. Consequently none of these actual
perfect tensors is weakly point-equivariantly
semiregular. This includes arbitrary actual
differentials and Postnikov couplings in $N$.

\paragraph{1. Actual derived duality and its signs.}
Set
\[
 N^\vee=R\mathcal Hom_P(N,\mathcal O_P).
\]
Perfect duality gives the global isomorphism
$N\otimes^{\mathbf L}N^\vee
 \simeq R\mathcal Hom_P(N,N)$.
The actual identity of $N$ defines
\[
 c_N:\mathcal O_P\longrightarrow
                       N\otimes^{\mathbf L}N^\vee,
 \qquad
 \epsilon_N:N^\vee\otimes^{\mathbf L}N
                           \longrightarrow\mathcal O_P.
 \tag{F1S32.4}
\]
The second map is evaluation. Let
$\widetilde\epsilon_N
=\epsilon_N\circ\tau_{N,N^\vee}$, where $\tau$
is the Koszul symmetry.
For an actual $\alpha:F\to F[2]$ define
\[
\begin{split}
 \operatorname{ptr}_N(\alpha):\quad
 E&\xrightarrow{1_E\otimes c_N}
               E\otimes^{\mathbf L}N\otimes^{\mathbf L}N^\vee\\
  &\xrightarrow{\alpha\otimes1_{N^\vee}}
               E[2]\otimes^{\mathbf L}N\otimes^{\mathbf L}N^\vee
   \xrightarrow{1_{E[2]}\otimes\widetilde\epsilon_N}E[2].
\end{split}
 \tag{F1S32.5}
\]
This uses the given object $N$ and its actual dual.
The identity-to-internal-Hom construction makes
$c_N$ canonical under quasi-isomorphism, so these
are globally specified derived morphisms.

The signs can be checked on any bounded local
complex $V^\bullet$ of finite free modules
representing $N$. Its dual has
\[
 (V^\vee)^j=(V^{-j})^\vee,\qquad
 d_{V^\vee}^{\,j}=(-1)^{j+1}(d_V^{-j-1})^\vee.
 \tag{F1S32.6}
\]
For a local basis $e_{a,u}$ in degree $a$, coevaluation
sends $1$ to $\sum_{a,u}e_{a,u}\otimes e_{a,u}^\vee$.
The two contributions to its differential cancel
by (F1S32.6), including the sign $(-1)^a$ in the
tensor differential. Evaluation
$V^\vee\otimes V\to\mathcal O_P$ is the usual
degreewise pairing.
Right evaluation after the symmetry instead gives
\[
 \widetilde\epsilon_N(e_{a,u}\otimes e_{a,v}^\vee)
                      =(-1)^a\delta_{uv}.
 \tag{F1S32.7}
\]
The two zigzag identities for (F1S32.4) use the
left evaluation and are the degreewise dual-basis
identities; hence they are actual identities of
complexes before passage to the derived category.
It follows, in particular, that the categorical trace is
\[
             \operatorname{tr}_N(1_N)
                  =\sum_a(-1)^a\operatorname{rk}V^a=r.
 \tag{F1S32.8}
\]
Thus the relevant coefficient is the signed rank,
including for shifted complexes.

For $\xi\in\operatorname{Ext}^2_P(E,E)$ and
$v\in\operatorname{Ext}^2_P(N,N)$ the composites give
\[
 \operatorname{ptr}_N(\xi\otimes1_N)=r\xi,\qquad
 \operatorname{ptr}_N(1_E\otimes v)
       =\operatorname{ev}_E\bigl(\operatorname{tr}_N(v)\bigr),
 \quad
 \operatorname{tr}_N(v)\in H^2(P,\mathcal O_P).
 \tag{F1S32.9}
\]
Here the trace of $v$ is explicitly
$\widetilde\epsilon_N[2]\,(v\otimes1_{N^\vee})c_N$.
The total degree is two, so moving its shift through
either tensor factor contributes $(-1)^{2a}=1$.
The equalities hold for the global composites
(F1S32.4)--(F1S32.5); they do not require the
cohomology sheaves of $N$ to split.

\paragraph{2. Tensor Atiyah classes on the actual complexes.}
Write the first-jet class as a degree-zero derived map
\[
 \operatorname{At}_{\rm jet}(M):
                 M\longrightarrow M\otimes^{\mathbf L}\Omega_P^1[1].
\]
With the input/output diagonal convention of S24, the
actual tensor identity is
\[
\begin{split}
 \operatorname{At}_{\rm jet}(E\otimes^{\mathbf L}N)
   ={}&(1_E\otimes\tau_{\Omega_P^1[1],N})
                 (\operatorname{At}_{\rm jet}(E)\otimes1_N)\\
      &+1_E\otimes\operatorname{At}_{\rm jet}(N).
\end{split}
 \tag{F1S32.10}
\]
The symmetry explicitly places the shifted form
factor last. To obtain this identity on the
given derived tensor, use locally free resolutions
of the two complexes with their actual differentials.
On each local frame the tensor connection is
$\nabla_E\otimes1+1\otimes\nabla_N$.
Its graded commutator with the tensor differential
is the sum of the two corresponding commutators.
Differences of the local connections on overlaps
satisfy the same sum. These commutators and
overlap differences give the first-jet extension
cocycle on the total resolution.
The formula is compatible with the differential
and transition maps. Its global extension comparison
can also be specified without a choice of connections.
Write $J(M)$ for first jets and $K_M=M\otimes\Omega_P^1$.
On flat resolutions the actual quotient
\[
 Q_{E,N}=
 \bigl(J(E)\otimes J(N)\bigr)\big/
                      \bigl(K_E\otimes K_N\bigr)
\]
is an extension of $E\otimes N$ by
$(K_E\otimes N)\oplus(E\otimes K_N)$.
The two extension components are the tensors
of the given jet extensions.
Multiplication in the first diagonal neighborhood
defines $Q_{E,N}\to J(E\otimes N)$: the product
of the two diagonal ideals vanishes modulo
$I_\Delta^2$.
On kernels this morphism is the sum of the
two differential-form contributions.
Pushing the displayed extension out by that sum
therefore gives precisely the Baer sum represented
by (F1S32.10), with its written shift symmetry.
These are global maps of the jet complexes;
flat resolutions only compute their derived tensors.
Both descriptions retain the actual differentials.

Represent $\kappa$ as the derived map
$\widetilde\kappa:\Omega_P^1[1]\to\mathcal O_P[2]$.
Naturality of the symmetry in (F1S32.10) replaces
$\tau_{\Omega_P^1[1],N}$, after contraction, by
$\tau_{\mathcal O_P[2],N}$.
Its sign is $(-1)^{2a}=1$ on degree $a$.
Using the corrected dictionary
$\operatorname{ev}_M(\kappa)
=-\widetilde\kappa\operatorname{At}_{\rm jet}(M)$
therefore gives
\[
 \operatorname{ev}_F(\kappa)
   =\operatorname{ev}_E(\kappa)\otimes1_N
                 +1_E\otimes\operatorname{ev}_N(\kappa).
 \tag{F1S32.11}
\]
This explicitly accounts for the shift which turns
the first-jet degree into the even degree of the
ambient obstruction.

Put
\[
 \tau_N(\kappa)=
          \operatorname{tr}_N\bigl(\operatorname{ev}_N(\kappa)\bigr)
                             \in H^2(P,\mathcal O_P).
\]
The scalar class $\eta$ acts once on $F$ by
multiplication. Equations (F1S32.9)--(F1S32.11) give
\[
 \operatorname{ptr}_N\operatorname{ev}_F(\kappa+\eta)
   =-r\,\operatorname{ob}_E(\kappa)
                +\operatorname{ev}_E\bigl(\tau_N(\kappa)+r\eta\bigr).
 \tag{F1S32.12}
\]
No identification of $\tau_N$ with a determinant
class is needed. S31's actual projector kills its
scalar evaluation and preserves
$\operatorname{ob}_E(\kappa)$ by naturality of the
quotient unit $\mathcal O_P\to E$. Applying that
projector proves (F1S32.2) for every perfect $N$.

\paragraph{3. The actual degree-four component forces nonzero virtual rank.}
S14 computed
\[
 \operatorname{ch}_j(E)=0\quad(j<4),\qquad
 \operatorname{ch}_4(E)=w+\frac{\beta^4}{72},
 \qquad 0\ne w\in W_K(P),\qquad \beta w=0.
 \tag{F1S32.13}
\]
Multiplicativity of the Chern character for the
actual derived tensor gives
\[
 \bigl(e^B\operatorname{ch}(F)\bigr)_4
               =r\left(w+\frac{\beta^4}{72}\right).
 \tag{F1S32.14}
\]
Let $q_4$ be the coefficient of $x^4$ in $Q(x)$.
Taking the same component of (F1S32.3) and then
applying the actual cup-product map
$\beta\wedge:H^8(P,\mathbf Q)\to H^{10}(P,\mathbf Q)$
gives
\[
          \frac r{72}\,\beta^5=q_4\beta^5,
 \tag{F1S32.15}
\]
because $\beta t=0$ for the determinant classes.
The class $\beta^5$ is nonzero: otherwise
$\beta^8=\beta^5\beta^3$ would be zero, contradicting
the positive polarization volume.
Hence $q_4=r/72$ and the unwedged equality yields
\[
                         \lambda t=rw.
 \tag{F1S32.16}
\]
Its left side is nonzero, so $r\ne0$.
This comparison uses a specified cohomological map
and the actual entire character requirement.

\paragraph{4. A full trace-kernel image with its actual surjection.}
For the given $B$, use the same linear map as S31,
\[
 a_B(\kappa)=\kappa+\kappa\mathbin{\lrcorner}B
 \in H^1(P,T_P)\oplus H^2(P,\mathcal O_P)\subset HH^2(P).
\]
The actual polarized Weil tangent annihilates
$\beta$ and both determinant classes.
Writing $w_F=\lambda t+Q(\beta)$, the finite
exponential identity gives
\[
 a_B(\kappa)\mathbin{\lrcorner}\operatorname{ch}(F)
  =\bigl(D_\kappa+(D_\kappa B)\wedge\bigr)
                      (e^{-B}w_F)=0.
 \tag{F1S32.17}
\]
BF's full characteristic--trace identity holds for
perfect complexes, so
\[
 U_F:=\operatorname{ev}_F(a_B(\mathfrak D_P))
                  \subseteq\ker\sigma_F^2.
\]
The map
\[
 \left.-\frac1r\,p_{\rm ns}\operatorname{ptr}_N\right|_{U_F}:
 U_F\twoheadrightarrow W_E,\qquad
 \operatorname{ev}_F(a_B(\kappa))
                         \longmapsto\operatorname{ob}_E(\kappa)
 \tag{F1S32.18}
\]
is the restriction of a specified linear map on
the whole self-Ext group.
It is surjective by the definition of $W_E$ and
(F1S32.2), proving $\dim U_F\ge16-r_C$.
No vanishing on the kernel of the original
period map is required to define this surjection.

Every point functor $t_x^*(-)\otimes M$,
$M\in\operatorname{Pic}^0(P)$, acts trivially on
$HH^2(P)$ by the invariant-form model and
$\operatorname{At}(M)=0$.
Naturality of the characteristic map therefore
fixes every element of $U_F$ under any specified
coherent lift on $F$.
It follows that
\[
 \dim\bigl(\ker\sigma_F^2\cap
                      \operatorname{Ext}^2_P(F,F)^G\bigr)
                              \ge16-r_C\ge4.
 \tag{F1S32.19}
\]
The R05 identity-sector and induction/forgetting
comparison identifies this with the weak
equivariant trace kernel. The entire argument
applies to the given perfect complex, irrespective
of its negative self-Ext; a candidate meeting the
gluability requirement still has the displayed
nonzero trace kernel.

\paragraph{5. Concrete complexes with all actual couplings retained.}
For a fixed rational divisor $D$, let each term
$V^a$ of an actual bounded complex be a finite
sum of isogeny line pushforwards $h_*L$ with
$h^*D=c_1(L)$, or one of the vector bundles
with the actual common-$D$ filtration of S31.
Give this complex any actual differential with
square zero.
S31 proves $\operatorname{ch}(V^a)
=\operatorname{rk}(V^a)e^D$, and therefore
\[
 \operatorname{ch}(N)
       =\sum_a(-1)^a\operatorname{ch}(V^a)
       =r e^D,\qquad
 e^{-D}\operatorname{ch}(E\otimes^{\mathbf L}N)
       =r\operatorname{ch}(E).
 \tag{F1S32.20}
\]
For $r\ne0$ this is the full S14 character with
$B=-D$, so (F1S32.19) applies.
The same character identity holds for any actual
finite iterated cone, finite direct sum or shift
formed from these common-$D$ bundles, by the
specified triangle relations in $K_0(P)$.
The maps (F1S32.4)--(F1S32.12) are still formed on
that actual resulting object; all its coupling
morphisms remain present.
Here $K_0$ is used only to evaluate its Chern
character. The obstruction comparison itself is
the derived duality composite on the object.

\paragraph{Primary-source and framework scope.}
The duality maps and signs are the ones constructed
in the Stacks Project, Lemma~20.50.5 (08DQ),
Example~20.50.7 (0FPC) and Lemma~15.74.2 (0FNK);
the proof of the lemma, the example's actual maps
and the full graded-sign proof were freshly read.
The example omits its zigzag verification; the
degreewise verification is given above.
The cup/composition compatibility was also
checked against the complete proof of Lemma~20.51.2
(0FVB).

The BF input and its corrected diagonal/HKR sign
are precisely those audited in S31:
Proposition~4.1.7, Theorem~5.1.3,
Corollary~5.2.3 and the full characteristic--trace
identity in Section~6 of
\emph{The global decomposition theorem for Hochschild
(co-)homology of singular spaces via the
Atiyah--Chern character}, Adv.\ Math.~217 (2008),
243--281, DOI 10.1016/j.aim.2007.06.013.
The tensor identity is verified here on the
actual complexes with the shifted form factor
and every symmetry displayed.
The quotient-unit projector, normal image,
whole character and weak comparison are the
specified S18/S24/S29/S14/R05 maps.

